\appendix

\section{Factors product vs Tensor product}
\label{app:cost}\label{app:section2}

%Tractability and efficiency of BNs rely on the properties of the operations on factors. In particular, the definition of factors product (that we denote $\odot$) \emph{inherently shares variables}.

We stress how much the product of factors in the theory of Bayesian networks (noted $\odot$ in our paper) differs from the tensor product of matrices. 
Consider the following BN (corresponding to the term in~\cite[Example 9.7]{popl24}), whose semantics is $\Pr(X_1, X_2,Y_1,Y_2,Y_3)$, a table with $2^5$ entries.

\begin{adjustbox}{width=0.7\linewidth, center=\linewidth}
\begin{tikzpicture}
\begin{genscope}
\node (Y1) at (2,0) {$Y_1$};
\node (Y2) at (6,0) {$Y_2$};
\node (Y3) at (10,0) {$Y_3$};
\node (X1) at (4,-2) {$X_1$};
\node (X2) at (8,-2) {$X_2$};
\draw[->] (Y1) edge (X1);
\draw[->] (Y1) edge (X2);
\draw[->] (Y2) edge (X1);
\draw[->] (Y2) edge (X2);
\draw[->] (Y3) edge (X1);
\draw[->] (Y3) edge (X2);
\end{genscope}
\end{tikzpicture}
\end{adjustbox}
The semantics $\den{X_1}$ of the node $X_1$ is $\Pr(X_1 \mid Y_1,Y_2,Y_3)$ (a factor $\phi_1(x_1,y_1,y_2,y_3)$), which can be seen as a stochastic matrix of size $2^4$.
Similarly, the semantics $\den{X_2}$ of the node $X_2$ is $\Pr(X_2 \mid Y_1,Y_2,Y_3)$ and can be seen as a stochastic matrix with $2^4$ entries.

By using the respective definitions, it is immediate to check that:
\begin{itemize}
\item
The tensor product $\den {X_1} \otimes \den {X_2} $ yields a table with $2^8$ entries.
\item
The factors product $\den {X_1} \odot \den {X_2} $ yields a table with $2^5$ entries (because non compatible instantiations are never computed).
\end{itemize}

\subparagraph{On the computational cost of the semantical interpretation: $\odot$ vs $\otimes$.}

%BNs have a dual nature, as tools for \emph{efficient probabilistic inference} and as causal models.
\emph{Inference} algorithms on BNs aim at computing the semantics of the model through intermediate, partial computations, \emph{without ever computing the full joint probability} (as unfeasible in practice).

The approaches based on monoidal categories typically focus on \emph{causality}, not on the cost of inference.
A central role in the semantics is here played by a product $\otimes$ behaving as the tensor product of matrices.
Computing the semantics of $n$ binary random variables easily leads to intermediate computations of size much larger than $2^n$, the size of the full joint distribution.
Looking again at the BN above, consider the sub-net containing only the nodes $X_1$ and $X_2$.
Its factor-based semantics is $\den{X_1} \odot \den{X_2}$.
In a categorical setting, one computes $\den{X_1} \otimes \den{X_2}$, which is \emph{larger} than the full joint probability distribution (the worst-case scenario in BN's inference).
The point is that \emph{inference is not the focus} there.

%In this work, we focus on the inference , aiming for quantum Bayesian networks (QBNs) that integrate the factor-based semantics of BNs.
%Inference algorithms on BNs aim at computing the semantics of the model through intermediate, partial computations, without ever computing the full joint probability. %(whose size for a BN with "only" $64$ binary random variables is $2^{64}$)

\section{Proof of Proposition~\texorpdfstring{\ref{lem:qfactors_op}}{18}}
\label{app:section3}

\qfactors*
\begin{proof}
We only prove closeness, the rest being simple computations for each claimed equation.
\begin{itemize}
\item
Consider $\someqfact_1\compfact \someqfact_2$ with $\someqfact_i$ a \qfactor\ over $(\bX_i,\qreg_i)$.
We will use the completely positive map $\qcup[\L(\HH_{\qreg})]:\L(\HH_{\qreg})\tens\L(\HH_{\qreg})\to\C$ defined on the canonical orthonormal basis by
\begin{equation*}
\qcup[\L(\HH_{\qreg})](\be \tens \be') \defeq 1 \text{ if } \be=\be' \text{ and } 0 \text{ otherwise}
\end{equation*}
and extended to the general case by linearity.
Fix $\bx\in\setofrv{\bX_1\cup\bX_2}$, with $\bx_i$ its restriction to $\setofrv{\bX_i}$.
Then, $\someqfact_1(\bx_1)$ is a positive matrix in $\L(\HH_{\qreg_1})\cong\L(\HH_{\qreg_1\difset\qreg_2})\tens\L(\HH_{\qreg_1\cap\qreg_2})$ and $\someqfact_2(\bx_2)$ is a positive matrix in $\L(\HH_{\qreg_2})\cong\L(\HH_{\qreg_1\cap\qreg_2})\tens\L(\HH_{\qreg_2\difset\qreg_1})$.
We want positivity of $(\someqfact_1 \compfact \someqfact_2) (\bx) \in \L(\HH_{\qreg_1\symdif\qreg_2})\cong\L(\HH_{\qreg_1\difset\qreg_2})\tens\L(\HH_{\qreg_2\difset\qreg_1})$.
We remark that:
\begin{equation*}
(\someqfact_1 \compfact \someqfact_2) (\bx)
=
\left(\id[\L(\HH_{\qreg_1\difset\qreg_2})] \tens \qcup[\L(\HH_{\qreg_1\cap\qreg_2})] \tens \id[\L(\HH_{\qreg_2\difset\qreg_1})]\right)
\left(\someqfact_1(\bx_1) \tens \someqfact_2(\bx_2)\right)
\end{equation*}
As $\someqfact_1(\bx_1) \tens \someqfact_2(\bx_2)$ is positive and $\qcup$ is completely positive, $(\someqfact_1 \compfact \someqfact_2) (\bx)$ is positive.
\item
Consider a \qfactor\ $\someqfact$ over $(\bY\uplus\{X\},\qreg)$.
Fixing $\by\in\setofrv{\bY}$, as each $\someqfact(\by,x)$ is a positive matrix, it follows that their sum $(\sum_X \someqfact) (\by) = \sum_{x\in\setofrv{X}}\someqfact(\by,x)$ is positive too.
\item
Consider a \qfactor\ $\someqfact$ over $(\bX,\qreg\uplus\qreg')$.
For any $\bx\in\setofrv{\bX}$, $(\sum_{\qreg'} \someqfact) (\bx) = \ptr{\HH_{\qreg'}}\left(\someqfact(\bx)\right)$ is positive since $\someqfact(\bx)$ is positive and the (partial) trace is completely positive.
\qedhere
\end{itemize}
\end{proof}

\section{Around \QCPT s}
\label{app:equivalence_qbns}

Notice that our semantics of quantum Bayesian networks does not distinguish inputs and outputs of nodes, similarly to the factors semantics or to the relational models of Linear Logic.
For instance, the interpretation of a node with a quantum input $\qreg_1$, a quantum output $\qreg_2$ and a classical output $X$ is a function from $\setofrv{X}$ to positive matrices in $\Pos{\HH_{\qreg_1} \tens \HH_{\qreg_2}}$.

\subsection{Intuitions for \QCPT s}
\label{sec:QCPT}

We give some intuition for \cref{def:condition_qfactor_qbn}.
It is straightforward to associate a \qfactor\ to a quantum instrument (we do so in \cref{app:equivalence_qbns_proof}).
The reverse is not always possible, and \qCPTs are those \qfactor s corresponding to quantum instruments.

Consider a \qfactor\ $\someqfact$ over $(\{X\}\uplus\bY,\qreg)$.
We want it to be the interpretation of a node $X$ in a quantum Bayesian network: given a density matrix $\M\in\D(\HH_\qreg)$ and a value $\by\in\setofrv{\bY}$, we want a probability distribution on $X$.
A natural way to do so is by seeing $\M$ as a \qfactor\ over $(\emptyset,\qreg)$ (\ie\ the function sending the unique element of a singleton to $\M$): then, $\someqfact\compfact\M$ is a \qfactor\ over $(\{X\}\cup\bY,\emptyset)$, \ie\ each $(\someqfact\compfact\M)(x,\by)$ is in $\nnReal$.
Therefore, we say that $\someqfact$ is a \qCPT\ for $X$ given $(\bY, \qreg)$ if for all $\M\in\D(\HH_\qreg)$ and $\by\in\setofrv{\bY}$, $\sum_{x\in\setofrv{X}}(\someqfact\compfact\M)(x,\by) = 1$.

Similarly, consider a \qfactor\ $\someqfact$ over $(\bX,\qreg\uplus\qreg')$ that we want to associate to a node $\qreg$ of a network.
Hence, for every density matrix $\M\in\D(\HH_{\qreg'})$ and value $\bx\in\setofrv{\bX}$, we should get a density matrix in $\D(\HH_{\qreg})$.
As before, seeing $\M$ as a \qfactor\ over $(\emptyset,\qreg')$, let us look at the \qfactor\ $\someqfact\compfact\M$ over $(\bX,\qreg)$.
For $\someqfact$ to be a \qCPT\ for $\qreg$ given $(\bX,\qreg')$, we require that for all $\M\in\L(\HH_{\qreg'})$ and $\bx\in\setofrv{\bX}$, the positive matrix $(\someqfact\compfact\M)(\bx)$ has trace $1$.

\Cref{def:condition_qfactor_qbn} is simply a reformulation of the two conditions found above.

\begin{remark}
While complete positivity is relevant for instruments, for \qfactor s we only need to consider positive matrices.
This is because given $\Eop:\L(\C)\to\L(\HH)$---which is isomorphic to a matrix in $\L(\HH)$---complete positivity collapses to positivity.
% (as eigenvalues of a tensor product of matrices are the products of the eigenvalues of these two matrices).
\end{remark}

\subsection{Equivalence between Instrument- and \Qfactor-based \qqbn s}
\label{app:equivalence_qbns_proof}

We prove here a claim from \cref{sec:qbn}: our definition of quantum Bayesian networks with \qfactor s is equivalent to the one with quantum instruments from~\cite{GBN2014}, sketched in \cref{sec:qbn_instrument}.

\begin{proposition}
\label{lem:bij_2_QBNs}
There is a one-to-one correspondence between instrument- and \qfactor-based Bayesian networks, which preserves the semantics.
\end{proposition}

As explained in \cref{sec:choi_iso}, the two syntaxes consider the same underlying DAGs (up to inconsequential details).
What we have to do to obtain \Cref{lem:bij_2_QBNs} is:
\begin{itemize}
\item
Giving a bijection $\Psi$ between a \qCPT associated to a node in our presentation, and instruments associated to the same node in~\cite{GBN2014}.
\item
Proving that composition of instruments corresponds to product of their images by $\Psi$.
\end{itemize}

This bijection $\Psi$ is a simple adaptation of the Choi-Jamiolkowski isomorphism to a setting with classical inputs.
This builds upon the fact that \qCPTs are defined similarly to Choi states.
We first define $\Psi$, for nodes associated to a \rv\ and nodes associated to a register.

\begin{lemma}
\label{lem:qcpt_qi}
\hfill
\begin{itemize}
\item
Let $X$ be a \rv, $\bY$ a set of \rvs, and $\qreg$ a register.
There is a bijection between families of quantum instruments $\{\Eop_\by \mid \by\in\setofrv{\bY}\}$ where each $\Eop_\by$ is an instrument $\{\Eop_{x,\by} \mid x\in\setofrv{X}\}$ from $\L(\HH_{\qreg'})$ to $\C$, and \qCPTs on $X$ given $(\bY,\qreg')$.
\item
Let $\bX$ be a set of \rvs, and $\qreg$ and $\qreg'$ two registers.
There is a bijection between families of quantum operations $\Eop_\bx$ from $\L(\HH_{\qreg'})$ to $\L(\HH_\qreg)$, and \qCPTs on $\qreg$ given $(\bX,\qreg')$.
\end{itemize}
\end{lemma}
\begin{proof}
We generalize the definition of \qCPTs (\cref{def:condition_qfactor_qbn}) to get both results as particular cases.
A \qfactor\ $\someqfact$ over $(\bX\uplus\bY,\qreg\uplus\qreg')$ is a \textdef{\qCPT\ for $(\bX,\qreg)$ given $(\bY, \qreg')$} if $\forall \by\in\setofrv{\bY}$:
\begin{equation*}
\left(\sum_{\bX,\qreg} \someqfact\right)(\by)
=
\id[\L(\HH_{\qreg'})]
\end{equation*}
%This generalizes the two cases in \cref{def:condition_qfactor_qbn}, and is equivalent to $\forall \by\in\setofrv{\bY}$, $\forall \be'\in\baseof{\qreg'}$:
%\begin{equation*}
%\sum_{\bx\in\setofrv{\bX}}\sum_{\be\in\baseof{\qreg}} \qfactfunc{\someqfact}(\bx,\by,\be,\be')\tr{\be}
%=
%\tr{\be'}
%\end{equation*}

We give a bijection $\Psi$ (along with its inverse $\Psi^{-1}$) between families of quantum instruments $\{\Eop_\by \mid \by\in\setofrv{\bY}\}$, where each $\Eop_\by$ is a quantum instrument $\{\Eop_{\bx,\by} \mid \bx\in\setofrv{\bX}\}$ from $\L(\HH_{\qreg'})$ to $L(\HH_\qreg)$, and \qCPTs for $(\bX,\qreg)$ given $(\bY, \qreg')$.
To simplify notations, we write $(\be_i)_{i}$ (\resp\ $(\be'_k)_{k}$) the canonical orthonormal basis of $\L(\HH_{\qreg})$ (\resp\ $\L(\HH_{\qreg'})$).
% Also, we note $\cdot$ the scalar product on these bases, \eg\ $\be_i \cdot \be_j = 1 \text{ if } i=j \text{ and } 0 \text{ otherwise}$.
We will use the following.
\begin{itemize}
\item
Call $\qcup[\L(\HH_{\qreg'})]:\L(\HH_{\qreg'})\tens\L(\HH_{\qreg'})\to\C$ the completely positive map defined by
\begin{equation*}
\qcup[\L(\HH_{\qreg'})](\be'_k \tens \be'_l) \defeq 1 \text{ if } k=l \text{ and } 0 \text{ otherwise}
\end{equation*}
and extended to the general case by linearity.
\item
Set $\qcap[\L(\HH_{\qreg'})]\defeq \sum_k \be'_k\tens \be'_k$, which is a positive matrix in $\L(\HH_{\qreg'})\tens\L(\HH_{\qreg'})$.
\end{itemize}

Let us now define $\Psi$ and $\Psi^{-1}$.
\begin{itemize}
\item
Given a family of quantum instruments $\{\Eop_\by \mid \by\in\setofrv{\bY}\}$ from $\L(\HH_{\qreg'})$ to $L(\HH_\qreg)$, we set $\Psi(\{\Eop_{\bx,\by} \mid \bx\in\setofrv{\bX}\})$ the \qfactor\ over $(\bX\uplus\bY,\qreg\uplus\qreg')$ defined by:
\begin{equation*}
\left(\Psi(\{\Eop_{\bx,\by} \mid \bx\in\setofrv{\bX}\})\right)(\bx,\by)
\defeq
(\id[\L(\HH_{\qreg'})] \tens \Eop_{\bx,\by})(\qcap[\L(\HH_{\qreg'})])
\end{equation*}
%Let us denote by $\someqfact$ the \qfactor\ $\Psi(\{\Eop_{\bx,\by} \mid \bx\in\setofrv{\bX}\})$.
%These $\someqfact(\bx,\by)$ are positive, as $\qcap$ is positive and each $\Eop_{\bx,\by}$ is completely positive.
%We prove they are \qCPT s for $(\bX,\qreg)$ given $(\bY,\qreg')$, \ie\ that $\sum_{\bx,i}\qfactfunc{\someqfact}(\bx,\by,\be_i,\be'_k)\tr{\be_i} = \tr{\be'_k}$.
%\begin{flalign*}
%\sum_{\bx,i}\qfactfunc{\someqfact}(\bx,\by,\be_i,\be'_k)\tr{\be_i}
%&=
%\sum_{\bx,i}(\id[\L(\HH_{\qreg'})] \tens \Eop_{\bx,\by})(\qcap[\L(\HH_{\qreg'})]) \cdot (\be'_k\tens\be_i)\tr{\be_i}
%\\
%&=
%\sum_{\bx,i,l} (\id[\L(\HH_{\qreg'})] \tens \Eop_{\bx,\by})(\be'_l\tens \be'_l) \cdot (\be'_k\tens\be_i)\tr{\be_i}
%\\
%&=
%\sum_{\bx,i,l} (\be'_l \tens \Eop_{\bx,\by}(\be'_l)) \cdot (\be'_k\tens\be_i)\tr{\be_i}
%\\
%&=
%\sum_{\bx,i} \Eop_{\bx,\by}(\be'_k) \cdot \be_i\tr{\be_i}
%\\
%&=
%\sum_{\bx}\tr{\Eop_{\bx,\by}(\be'_k)}
%=
%\tr{\sum_{\bx}\Eop_{\bx,\by}(\be'_k)}
%\\
%&=
%\tr{\be'_k}
%\tag*{as $\sum_\bx \Eop_{\bx,\by}$ is trace-preserving}
%\end{flalign*}

\item
Given a \qCPT\ $\someqfact$ for $(\bX,\qreg)$ given $(\bY, \qreg')$, we define for each $\by\in\setofrv{\bY}$ the following set of maps:
\begin{equation*}
\Psi^{-1}_{\by}(\someqfact)
\defeq
\{ \M \mapsto (\qcup[\L(\HH_{\qreg'})] \tens \id[\L(\HH_\qreg)]) (\M \tens \someqfact(\bx,\by)) \mid \bx\in\setofrv{\bX}\}
\end{equation*}
and pose $\Psi^{-1}(\someqfact) \defeq \{\Psi^{-1}_{\by}(\someqfact)\mid \by\in\setofrv{\by}\}$.
%Then, $\Psi^{-1}_{\by}(\someqfact)$ is a set of linear maps from $\L(\HH_{\qreg'})$ to $\L(\HH_\qreg)$, each completely positive since $\qcup$ is completely positive and $\someqfact(\bx,\by)$ is a positive matrix.
%We prove that their sum is trace-preserving, hence $\Psi^{-1}_{\by}(\someqfact)$ is a quantum instrument.
%It suffices to prove it for the basis $(\be'_k)_k$ thanks to linearity.
%We have:
%\begin{flalign*}
%\sum_\bx \Psi^{-1}_{\by}(\someqfact)(\be'_k)
%&=
%\sum_\bx (\qcup[\L(\HH_{\qreg'})] \tens \id[\L(\HH_\qreg)])
%(\be'_k \tens (\sum_{i,l}\qfactfunc{\someqfact}(\bx,\by,\be_i,\be'_l) \be'_l \tens \be_i))
%\\
%&=
%\sum_{\bx,i,l}\qfactfunc{\someqfact}(\bx,\by,\be_i,\be'_l)
%(\qcup[\L(\HH_{\qreg'})] \tens \id[\L(\HH_\qreg)])
%(\be'_k \tens \be'_l \tens \be_i)
%\\
%&=
%\sum_{\bx,i}\qfactfunc{\someqfact}(\bx,\by,\be_i,\be'_k) \be_i
%\end{flalign*}
%Thence, $
%\tr{\sum_\bx\Psi^{-1}_{\by}(\someqfact)(\be'_k)}
%=
%\sum_{\bx,i}\qfactfunc{\someqfact}(\bx,\by,\be_i,\be'_k)\tr{\be_i}
%=
%\tr{\be'_k}$ as $\someqfact$ is a \qCPT.
\end{itemize}
By similar computations to that involved for the Choi-Jamiolkowski isomorphism, $\Psi$ sends quantum instruments to \qCPTs, $\Psi^{-1}$ sends \qCPTs to quantum instruments, and $\Psi^{-1}$ is the inverse of $\Psi$ (the later computation using the yanking equations of $\qcup$ and $\qcap$).
\end{proof}

Showing the composition of instruments results in the product of their images by $\Psi$ is also a simple computation (using the yanking equations of $\qcup$ and $\qcap$).
This proves \cref{lem:bij_2_QBNs}.

\begin{remark}
\Cref{lem:sem_qbn_dist} is a corollary of \cref{lem:bij_2_QBNs}, using that the semantics of an instrument-based Bayesian network is a probability distribution.
Another corollary is that \qCPTs are closed under product, because instruments are closed by composition.
\end{remark}

% We use the following intermediate result to prove \cref{lem:sem_qbn_dist}.
% \begin{lemma}
% Take $\mathcal{B}$ a sub-graph of a quantum Bayesian network which is closed under taking parents of its nodes.
% Then $\someqfact^{\mathcal{B}}$ is a \qCPT\ for $(\bX,\qreg)$ given $(\emptyset,\emptyset)$ where $\bX$ are the \rvs\ in $\mathcal{B}$ and $\qreg$ the quantum outputs of $\mathcal{B}$.
% \end{lemma}

% Idea from Fabio: putting $\tens$ as $\cup$ on objects and $\compfact$ on morphisms may give a SMCC describing factors!
